% !TEX root = ../notes.tex

\documentclass[../notes.tex]{subfiles}

\begin{document}

\section{March 11}
We began class by completing the proof of \Cref{thm:haar-measure}.

\subsection{Examples of Compact Groups}
To justify our study of compact groups, we should probably give some interesting examples which are not Lie groups.
\begin{definition}[profinite]
	Let $\{G_i\}$ be a countable collection of finite groups equipped with surjective homomorphisms $G_{i+1}\onto G_i$ for each $i$. Then we define the \textit{profinite} group $G$ to consist of sequences $(g_i)\in\prod_iG_i$ for which $g_i\in G_i$ is the image of $g_{i+1}$. Thus, by the Axiom of choice, there are canonical surjections $\pi_n\colon G\onto G_n$, and we may give $G$ a topology by declaring the kernels of each $\pi_n$ to be a neighborhood basis of the identity.
\end{definition}
\begin{remark}
	One may equivalently define the topology to be the subspace topology from the canonical injection $G\to\prod_iG_i$. It is not hard to check, then, that $G$ is closed in $\prod_iG_i$, so it follows that $G$ is compact. In general, any inverse limit of compact groups will produce a compact group.
\end{remark}
\begin{remark}
	One can check that $G^\circ=\{1\}$.
\end{remark}
\begin{remark}
	The group $G$ is infinite as long as $\lim_n\left|G_n\right|=\infty$. In fact, as soon as $\lim_n\left|G_n\right|=\infty$, one can check that $G$ is uncountable and has no isolated points.
\end{remark}
\begin{example}
	For any prime $p$, we can define the group $\ZZ_p$ to be given by the countable collection
	\[\cdots\to\ZZ/p^3\ZZ\to\ZZ/p^2\ZZ\to\ZZ/p\ZZ.\]
	In fact, $\ZZ_p$ is a ring, so for any algebraic group $G$ over $\ZZ$, one also has a profinite group $G(\ZZ_p)$.
\end{example}
\begin{example}
	The group $\op{Gal}(\overline\QQ/\QQ)$ can be constructed as the limit of the groups $\op{Gal}(K/\QQ)$, where $K/\QQ$ is a Galois extension. Thus, $\op{Gal}(\overline\QQ/\QQ)$ admits a natural structure of a profinite group.
\end{example}
Here are some more exotic examples.
\begin{example}
	The inverse limit of the groups $\{\RR/\ZZ\times\ZZ/p^\bullet\ZZ\}$ is $\RR/\ZZ\times\ZZ_p$, which is a compact group.
\end{example}
\begin{example}
	Let $G$ be the inverse limit $\lim\RR/p^\bullet\ZZ$, where the projections $\RR/p^{n+1}\ZZ\to\RR/p^n\ZZ$ are the canonical ones (induced by $\id\colon\RR\to\RR$). Then $G$ is an inverse limit of compact groups, so it is a compact group. We claim that $G$ lives in a short exact sequence
	\[0\to\ZZ_p\to G\to\RR/\ZZ\to0.\]
	Indeed, the map $G\to\RR/\ZZ$ is the canonical one, so it only remains to compute the kernel. Well, for each $n$, we see that the kernel of $\RR/p^n\ZZ\to\RR/\ZZ$ is represented by $\ZZ/p^n\ZZ$, so taking the limit gives $\ZZ_p$.
\end{example}

\subsection{Representation Theory of Compact Groups}
Our integration theory now has the following applications.
\begin{corollary}
	Fix a separable compact topological group $G$. Then the finite-dimensional representations of $G$ are unitary and completely reducible.
\end{corollary}
\begin{proof}
	Repeat the proof of \Cref{cor:reps-semisimple}.
\end{proof}
\begin{theorem}[Peter--Weyl] \label{thm:pw-compact}
	Fix a separable compact topological group $G$.
	\begin{listalph}
		\item The collection of finite-dimensional irreducible representations $\op{IrRep}G$ is countable, and
		\[L^2(G)=\widehat{\bigoplus_\rho}(V\otimes V^*).\]
		\item Set $L^2_{\mathrm{alg}}(G)\coloneqq\bigoplus_\rho(V\otimes V^*)$. Then $L^2_{\mathrm{alg}}(G)$ is dense in $C(G)$, where $C(G)$ has been given the supremum norm.
	\end{listalph}
\end{theorem}
\begin{proof}
	The same arguments for \Cref{thm:peter-weyl} and \Cref{cor:l2-alg-dense} apply, as soon as we have a continuous approximation of the identity. Well, let $\{U_i\}$ be a neighborhood basis of the identity; we may assume that $\overline U_{i+1}\subseteq U_{i}$ for each $i$ because $G$ is locally compact. Then Urysohn's lemma grants us functions $f_i\colon G\to[0,1]$ such that $\op{supp}f_i\subseteq U_i$ but $f_i|_{U_{i+1}}=1$. Thus, $\int_Gf(g)\,dg\ge0$, and we can rescale these functions to be our continuous approximation of the identity.
\end{proof}
\begin{remark}
	Suppose $G$ is profinite, then any continuous representation of $G$ maps to $\op{GL}(V)$ for some complex vector space $V$. A ``no small subgroups'' argument is able to show that this representation thus factors through a finite quotient of $G$.
\end{remark}
As an application, we may classify our compact groups!
\begin{corollary}
	Any compact group $G$ is an inverse limit of compact Lie groups.
\end{corollary}
\begin{proof}
	By \Cref{thm:pw-compact}, we may enumerate the irreducible representations of $G$ as $\{\rho_i\}_{i\in\NN}$. Let $K_m$ be the kernel of the representation $\bigoplus_{i\le m}\rho_i$, and define $G_m\coloneqq G/K_m$. We now combine two observations.
	\begin{itemize}
		\item Note that $G/K_m$ embeds as a closed subgroup of $\op U(\rho_1\oplus\cdots\oplus\rho_m)$, which turns out to be a Lie subgroup. Thus, $G_m$ is a compact Lie group.
		\item On the other hand, matrix coefficients of the $\rho_i$s separate points by \Cref{thm:pw-compact}. Indeed, it is enough to show that $L^2_{\mathrm{alg}}(G)$ separates points, but this is dense in $C(G)$, which certainly separates points (for example, by our construction of approximation to identity).
	\end{itemize}
	Combining the above two points shows us that the canonical map
	\[G\to\lim_{m\to\infty}G_m\]
	embeds $G$ into an inverse limit of compact Lie groups. We claim that this map is actually surjective, which will complete the proof. Indeed, choose a sequence $(g_m)$ in the target. For each $m$, we see that the pre-image of $g_m\in G/K_m$ is the closed subset $g_mK_m\subseteq G$. By construction, this family of compact sets $\{g_mK_m\}_m$ is a descending chain of nonempty compact sets, so the finite intersection property for compact groups implies that $\bigcap_mg_mK_m$ is nonempty. Any $g\in G$ in this intersection will map to the sequence $(g_m)$!
\end{proof}

\subsection{The Hydrogen Atom}
We now switch topics. Electrons in an atom turn out to not exactly behave as particles (with well-defined position at all times). Instead, they kind of behave like they are in a cloud around the atom, with some probability of being in any particular location. We are interested in calculating the probability distribution $p(x,y,z,t)$ of an electron living in a hydrogen atom, which satisfies
\[\int_{\RR^3}p(x,y,z,t)\,dx\,dy\,dz=1\]
for all $t$. Schr\"odinger's quantum mechanics claims that $p$ arises from a more natural complex-valued function $\psi(x,y,z,t)$ from which one recovers $p$ from $p=\left|\psi\right|^2$. To calculate $\psi$, one has the following.
\begin{theorem}[Schr\"odinger's equation]
	The wave function $\psi$ satisfies
	\[i\frac{\del\psi}{\del t}=H\psi,\]
	where $H$ is a quantum Hamiltonian $-\frac12\Delta-\frac1r$.
\end{theorem}
\begin{remark}
	This $H$ arises from quantizing the classical Hamiltonian $\frac{p^2}2-\frac1r$, where the momentum $p=p_x^2+p_y^2+p_z^2$ is ``quantized'' to $\Delta=\del_x^2+\del_y^2+\del_z^2$.
\end{remark}
\begin{remark}
	The previous discussion implies that our wave function $\psi$ is required to satisfy
	\[\int_{\RR^3}\left|\psi(x,y,z,t)\right|^2\,dx\,dy\,dz=1\]
	for all $t$.
\end{remark}
One can solve this equation using separation of variables because the right-hand side does not depend on time. As such, let's quickly review how separation of variables works.
\begin{exe}
	Suppose that we have a Hermitian operator $A$ on a complex finite-dimen\-sional vector space $V$, and we solve the equation
	\[i\frac{\del v}{\del t}=Av.\]
\end{exe}
\begin{proof}
	Well, $A$ admits a diagonalization with eigenbasis $\{v_j\}$ and real eigenvalues $\{\lambda_j\}$. Then we may write $v(t)=\sum_jf_j(t)v_j$, so we see that
	\[if_j'(t)=\lambda_jf_j(t),\]
	which implies that $f_j(t)=f_j(0)e^{-\lambda_jt}$. We thus see that
	\[v(t)=\sum_jf_j(0)e^{-i\lambda_jt}v_j,\]
	as required.
\end{proof}
We now attempt to solve Schr\"odinger's equation. Suppose that we could find a basis $\{\psi_n\}$ of $L^2(\RR^3)$ satisfying $H\psi_n=E_n\psi_n$ for each $n$, where each $E_n$ is real. Then the same argument shows that
\[\psi(x,y,z,t)=\sum_na_ne^{-iE_nt}\psi_n(x,y,z).\]
Thus, we reduce ourselves to actually finding these functions $\psi$ with $H\psi=E\psi$. Note that we are interested in solving ``bounded states'' for our electron, which amounts to requiring $\psi$ to stay bounded near the origin, which we can see is equivalent to $E<0$.

We would like to exploit the symmetry of our system to find these functions $\psi$. In particular, the form of $\Delta$ is invariant under rotations, so we should try to work in spherical coordinates $(r,\theta,\varphi)$. In particular, it turns out that
\[\Delta=\Delta_r+\frac1{r^2}\Delta_{\mathrm{sph}},\]
where
\[\Delta_r=\del_r^2+\frac2r\del_r,\]
and
\[\Delta_{\mathrm{sph}}=\frac1{\sin^2\varphi}\del_\theta^2+\frac1{\sin\varphi}\del_\varphi\circ\sin\varphi\circ\del_\varphi,\]
which is the Laplace--Beltrami operator. We are now hunting for functions $\psi$ for which
\[\del_r^2\psi+\frac2r\del_r\psi+\frac2r\psi+\frac1{r^2}\Delta_{\mathrm{sph}}\psi=-2E\psi.\]
We may now apply separation of variables: we may write $\psi$ as a function of $r\in\RR^+$ and a unit vector $u$, and then we will pretend that $\psi$ can be spanned by functions of the form $\psi(r,u)=f(r)\xi(u)$, where $\xi$ is an eigenvector for $\Delta_{\mathrm{sph}}$ with eigenvalue $-\lambda$. Then we are interested in solving the equations $\Delta_{\mathrm{sph}}\xi=-\lambda\xi$ and
\[f''(r)+\frac2rf'(r)+\left(\frac2r-\frac\lambda{r^2}+2E\right)f(r)=0.\]
We will solve these equations, using the representation theory of $\op{SO}(3)$, next class.

\end{document}